
Multi-modal generative modelling is a crucial area of research in machine learning that aims to develop models capable of generating data according to multiple modalities, such as images, text, audio, and more. This is important because real-world observations are often captured in various forms, and combining multiple modalities describing the same information can be an invaluable asset. For instance, images and text can provide complementary information in describing an object, audio and video can capture different aspects of a scene. Multi-modal generative models can also help in tasks such as data augmentation \citep{he2023is, azizi2023synthetic, sariyildiz2023fake}, missing modality imputation \citep{pmlr-v97-antelmi19a, 9680207, zhang2023unified, Tran_2017_CVPR}, and conditional generation \citep{huang2022eccv, lee2018harmonizing}.

Multi-modal models have flourished over the past years and have seen a tremendous interest from academia and industry, especially in the content creation sector. Whereas most recent approaches focus on specialization, by considering text as primary input to be associated mainly to images \citep{Rombach_2022_CVPR, saharia2022photorealistic, ramesh2022hierarchical, tao2022dfgan, Wu_2022_CVPR, nichol2022glide, chang2023muse} and videos \citep{blattmann2023align, hong2023cogvideo, singer2022makeavideo}, in this work we target an established literature whose scope is more general, and in which all modalities are considered equally important. A large body of work rely on extensions of the \gls{VAE} \citep{kingmaW13vae} to the multi-modal domain: initially interested in learning joint latent representation of multi-modal data, such works have mostly focused on generative modeling.
Multi-modal generative models aim at \textit{high-quality} data generation, as well as generative \textit{coherence} across all modalities. These objectives apply to both joint generation of new data, and to conditional generation of missing modalities, given a disjoint set of available modalities.


In short, multi-modal \glspl{VAE} rely on combinations of uni-modal \glspl{VAE}, and the design space consists mainly in the way the uni-modal latent variables are combined, to construct the joint posterior distribution. Early work such as \cite{mvaeWu2018} adopt a product of experts approach, whereas others \cite{mmvaeShi2019} consider a mixture of expert approach. Product-based models achieve high generative quality, but suffer in terms of both joint and conditional coherence. This was found to be due to experts mis-calibration issues \citep{mmvaeShi2019, sutter2021generalized}. On the other hand, mixture-based models produce coherent but qualitatively poor samples.
A first attempt to address the so called \textbf{coherence--quality tradeoff} \citep{Daunhawer2021} is represented by the mixture of product of experts approach \citep{sutter2021generalized}. However recent comparative studies \citep{Daunhawer2021} show that none of the existing approaches fulfill both the generative quality and  coherence criteria. A variety of techniques aim at finding a better operating point, such as contrastive learning techniques \citep{contrastShi2020}, hierarchical schemes \citep{Vasco2022}, total correlation based calibration of single modality encoders \citep{tcmvaehwang21}, or different training objectives \cite{jsdSutter2020}. More recently, the work in \citep{palumbo2023mmvae} considers explicitly separated shared and private latent spaces to overcome the aforementioned limitations.

By expanding on results presented in \citep{Daunhawer2021}, in \Cref{sec:motivation} we further investigate the tradeoff between generative coherence and quality, and argue that it is intrinsic to all variants of multi-modal \glspl{VAE}. We indicate two root causes of such problem: latent variable collapse \citep{alemi2018fixing, dieng2019avoiding} and information loss due to mixture sub-sampling. To tackle these issues, in this work, we propose in \Cref{sec:method} a new approach which uses a set of independent, uni-modal \textit{deterministic} auto-encoders whose latent variables are simply concatenated in a joint latent variable. Joint and conditional generative capabilities are provided by an additional model that learns a probability density associated to the joint latent variable.
We propose an extension of score-based diffusion models \citep{Song2020} to operate on the multi-modal latent space. We thus derive both forward and backward dynamics that are compatible with the multi-modal nature of the latent data. In \cref{sec:mld_var} we propose a novel method to train the multi-modal score network, such that it can both be used for joint and conditional generation. Our approach is based on a guidance mechanism, which we compare to alternatives. We label our approach \gls{MLD}.

Our experimental evaluation of \gls{MLD} in \Cref{sec:experiments} provides compelling evidence of the superiority of our approach for multi-modal generative modeling. We compare \gls{MLD} to a large variety of \gls{VAE}-based alternatives, on several real-life multi-modal data-sets, in terms of generative quality and both joint and conditional coherence. Our model outperforms alternatives in all possible scenarios, even those that are notoriously difficult because modalities might be only loosely correlated. 
% Although we are aware of concurrent work that explore the joint generation of multiple modalities \cite{Ruan2022, Hu2022}, also with diffusion models \citep{bao2023transformer}, such approaches are application specific, e.g. text-to-image, and essentially only target two modalities. Nevertheless, when possible, we compare our method to alternative approaches to multi-modal diffusion, and show superior performance of \gls{MLD}.
Note that recent work also explore the joint generation of multiple modalities \cite{Ruan2022, Hu2022}, but such approaches are application specific, e.g. text-to-image, and essentially only target two modalities. When relevant, we compare our method to additional recent alternatives to multi-modal diffusion \citep{bao2023transformer, wesego2023scorebased}, and show superior performance of \gls{MLD}.



